\chapter{Technical background}
\label{chap:background}

Detection tasks can be solved with a variety of solutions, each tailored to
specific problems and desired performance. The detection of chirps is no
exception. In this chapter, we will introduce some necessary background
information: How are fields of electric fish recorded and how does this relate
to chirp detection? What can we observe chirps if fish can freely interact? How
were chirps detected before? And what is an artificial neural network?

\section{How to record electric fish?}
\label{sec:recording}

Traditionally, electric fish recordings have utilized a differential setup with
two electrodes, one placed near the head and the other near the tail, to
capture the electric field generated by the fish. This setup, known as a
\enquote{chirp chamber}, is widely used in research (e.g.,
\cite{dunlapProductionAggressiveElectrocommunication2003,
fugereElectricSignalsSpecies2010}). Species like \textit{A.~leptorhynchus}
often remain in shelters when exposed to light, mimicking their natural
inactive phase, which reduces the need for physical restraint during
recordings. This method produces a nearly perfect, quasi-sinusoidal
\acrfull{EOD} from single, isolated individuals. Such a recording is ideally
centered around zero and noiseless enough so that every single period, i.e.,
every single discharge as well as the species specific waveform, is clearly
visible. However, it limits the study of social behaviors and interactions, as
animals cannot engage physically or electrically, making it less suitable for
research into communication.

To investigate such interactions, the traditional setup can be adapted by
housing two fish in a single tank divided by a mesh, allowing electrical but
not physical interaction (e.g., \cite{obotiWhyBrownGhost2023}). Electrodes are
placed in each compartment to record the electric fields. While these
recordings can capture the fish's instantaneous frequencies, they tend to be
noisier and are better analyzed using spectrograms rather than direct frequency
calculations. However, such setups are highly dependent on the experimental
conditions, including tank size and electrode placement, which can influence
the observed behaviors and may not accurately reflect natural interactions.
Smaller setups are often used to minimize electrode requirements, but this can
lead to unnatural behaviors not observable in the wild. Additionally, the
effectiveness of EOD recording varies with electrode spacing, and physical
barriers limit genuine interactions. A potential solution is to employ an
electrode grid, distributing many electrodes evenly throughout the setup to
enhance recording fidelity during interactions
\parencite{junLongtermBehavioralTracking2014, pedrajaUseSupervisedLearning2021,
henningerTrackingActivityPatterns2020,
madhavHighresolutionBehavioralMapping2018}.

Electrode grids, initially developed for field research, are equally effective
in controlled studies, where they ensure complete coverage of the recording
setup and allow free interactions between fish. These setups necessitate the
use of spectrograms for frequency analysis, which introduces challenges related
to spectrogram parameters and tracking individual fish EODs over time.
Solutions to these challenges have been proposed by
\textcite{henningerTrackingActivityPatterns2020} and
\textcite{madhavHighresolutionBehavioralMapping2018}, and further refined by
\textcite{raabAdvancesNoninvasiveTracking2022}, culminating in the
\href{https://github.com/tillraab/wavetracker}{\texttt{wavetracker}} software
package for automatic fundamental frequency tracking. This advancement
facilitated the first study on electrocommunication in unconstrained
interactions in \textit{A.~leptorhynchus} on a large dataset by
\textcite{raabElectrocommunicationSignalsIndicate2021}. However, chirp
detection still requires manual spectrogram analysis for fast frequency
changes.

As we transition to the next section, I will cover previous attempts to detect
chirps from the rich datasets generated by electrode grids. This exploration
will not only mark the challenges faced in earlier methodologies but allow a
discussion on the techniques that have emerged to overcome these obstacles. But
to understand the challenge of detecting chirps in the EOD of multiple fish, it
is crucial to first introduce how we can make chirps accessible in recordings
of multiple fish.

\section{How to observe chirps in naturalistic settings?} % (fold)
\label{sec:chirps}

In the \hyperref[chap:intro]{introduction} we already established that the main
features of the \acrshort{EOD} that are changed during a chirp are the
frequency of discharge as well as the amplitude. But this amplitude feature is
not easily accessible when multiple fish are recorded on the same electrode,
which usually happens when fish interact. In this case, the resulting signal is
a superposition of the EODs of all fish, where the amplitude is modulated by a
so-called \emph{beat}. The exact shape of this \acrfull{AM} is dependent on the
differences between the discharge frequencies of the fish, the distance of each
fish to the recording electrode and even the orientation as well as curvature
of the body of single fish. Particularly small and short amplitude modulations
during chirps often drown in this complexity (figure \ref{fig:multifish_chirps}
C) while the AM of larger chirps remains visible in some cases (figure
\ref{fig:multifish_chirps} G). But in these cases, the time scales mix with
those of body movements. So if such a modulation is the result of a chirp or
that of a brief attenuation by a non-conducting object, such as a stone, is not
always clear.

So, in most species the property of the signal that firstly changes the
\enquote{most} and secondly is not influenced by external factors is the
frequency modulation during a chirp. To approximate how frequency components
change over time, we can use spectrograms. To compute a spectrogram, a
\acrfull{STFT} is applied to the recorded EOD to estimate the frequencies and
their powers of a part of the signal. The main pitfall with this approach is
given by the uncertainty principle: To get a very precise estimate of the
powers of certain frequencies in a signal, we need to compute the STFT over a
sufficiently large chunk of the signal, as the duration of a single segment
sets the frequency resolution. As an example, including three seconds of a
signal into a single Fourier transform results in a very fine frequency
resolution of $1/3$ \si{\hertz}. That means that even fish that are just one
Hertz apart in their EOD$f$ would be separable from each other on a
spectrogram. But for each point in time, we would end up with just the average
of the frequency estimate in three seconds around this point. Given that a
chirp only lasts fractions of a second, it would be impossible to find chirps
on these spectrograms. On the other hand, if we resolve the time domain of a
spectrogram to be able to closely follow the frequency modulation of a chirp,
the frequency bands for individual fish would be so broad that they would merge
into each other. Now we see chirps well, but cannot tell which fish produced
them. The ideal spectrogram considers the optimal trade-off between the two: It
has sufficient temporal resolution to make chirps somewhat visible but also
maximizes the separability of fish. An example for such a spectrogram is
illustrated in figure \ref{fig:multifish_chirps} D for small chirps and figure
\ref{fig:multifish_chirps} H for large chirps. We see that particularly the
temporal evolution of small chirps become blurry as the temporal resolution is
decreased in favor to make fish separable. To resolve chirps while still
maintaining reasonable separability of individual fish, a spectrogram that uses
bins of approximately \SI{20}{\milli\second} corresponding to 4096 samples for
a signal with a \SI{20}{\kilo\hertz} sampling rate is ideal. This makes sense
given that chirps are usually not shorter than \SI{20}{\milli\second} and at
least a single frequency estimate per chirp is needed to resolve it well. The
second important factor is the overlap between the bins: If we estimate the
frequency once every \SI{20}{\milli\second} and just by chance place one bin at
the start and one bin at the end of the chirp, the frequency in these bins will
still be elevated compared to the baseline but the increase won't be that
pronounced. So ideally, the overlap of the windows in which the frequencies are
estimated is large and traded off against the availability of computational
power. A stronger overlap produces a finer spectrogram but requires more
computations and memory.

Now that we have introduced why chirps are interesting, which parameters define
chirps in naturalistic recording conditions and how we can make them
approachable by computing specific types of spectrograms, let us discuss
previous implementations of chirp detectors.

\section{How were chirps detected in the past?}
\label{sec:chirpdetectors}

Detecting chirps in the EOD of electric fish, particularly in a controlled
environment with restrained fish, has traditionally been straightforward. In
such settings, fish are placed in chirp chambers where their movements are
limited, ensuring stable EOD signals without overlap from other sources. The
most common technique involves band-pass filtering the EOD signal and
calculating the instantaneous frequency through zero-crossings, with chirps
identified by rapid frequency increases as illustrated in figure
\ref{fig:multifish_chirps} A and E.

\textcite{henningerStatisticsNaturalCommunication2018} developed a method to
detect chirps in the electric organ discharges of multiple fish on field
electrode grid recordings. They initially filtered each fish's EOD within a
$\pm$ \SI{7}{\hertz} range around the baseline frequency, and additionally
within a \SI{14}{\hertz} window positioned \SI{10}{\hertz} above the baseline
frequency to capture the chirp signals, which momentarily elevate into this
range. To identify chirps, they employed a \acrfull{RMS} filter to estimate the
signal's instantaneous amplitude or envelope, looking for peaks within this
\enquote{search envelope} that exceeded a predefined threshold. This method
also involved monitoring for concurrent drops in the EOD amplitude, with
simultaneous amplitude drops and search envelope peaks indicating chirps. Each
detected chirp was manually verified and attributed to the emitting fish based
on the assumption that such signal patterns uniquely occur during chirping.
This approach enabled the successful detection and assignment of chirps to
specific fish in complex multi-fish recordings during courtship behavior.

In their research on the role of chirps in \textit{A. leptorhynchus},
\textcite{obotiWhyBrownGhost2023} applied a technique akin to
\textcite{henningerStatisticsNaturalCommunication2018} for chirp detection,
focusing on identifying peaks in the power spectral density between the
fundamental frequency and the first harmonic of the EOD. Peaks surpassing a
noise threshold were classified as chirps, with each detection undergoing
manual verification. For chirp assignment to specific fish, Oboti
leveraged the recording setup's electrode spacing, ensuring one fish dominated
the signal on a particular electrode, thus facilitating accurate chirp
attribution. While effective in controlled lab settings, this method's
requirement for limited fish movement compromises its applicability to field
recordings and may not fully capture natural fish communication behaviors due
to the imposed movement restrictions.

Neither study discussed the false positive rates of their chirp detection
methods or the extent of manual verification needed. In field recordings, it is
possible to capture up to 25 fish simultaneously over several weeks, with fish
potentially chirping at rates of two to three times per second, sometimes for
hours. Henninger's method struggles with chirp assignment in recordings where
fish EOD frequencies are closely spaced. Oboti's approach, reliant on
restricting fish movement, probably faces challenges in accurately assigning
chirps, since chirps often occur when fish are in close proximity, making it
difficult to assign chirps based on the strongest electrode signal.
Additionally, both methods are susceptible to interference from electrical
noise, such as white noise bursts that mimic chirps and can interfere with the
EOD signal, complicating chirp identification without extensive manual review,
which is impractical for large datasets.

The following sections will detail the advancements in chirp detection and
different approaches that I applied in the preparation and in initial
experiments of this thesis.

\subsection{Dynamic filtering of the EOD for chirp detection}
\label{sec:filterdetector}

Our initial chirp detection approach, developed with Sina Prause, Alexander
Wendt, and Till Raab, utilized dynamic filtering based on simple filters and
parameters from \textcite{henningerStatisticsNaturalCommunication2018}. Unlike
Henninger's stationary search envelope, our method dynamically adjusted the
frequency range based on the presence of other fish, optimizing chirp detection
for each recording window. We also incorporated spatial information, focusing
on the electrode with the strongest signal for each fish, to enhance
\acrfull{SNR} and avoid frequency overlap with other fish. The detection relied
on peaks occurring simultaneously in three features: The dynamic search
envelope, the envelope filtered around the EOD baseline, and the instantaneous
frequency of the baseline EOD (this code is available on GitHub:
\url{https://github.com/weygoldt/gp-chirpdetector}).

This method was effective for simple scenarios with two fish with a wide gap in
their EOD$f$s. However, it struggled with fast moving fish due to noisy
baseline envelopes and unreliable spatial information, leading to a high false
positive rate and difficulty in assigning chirps to the correct fish,
especially for long chirps. Given these limitations and the reliance on
hand tuned parameters, I shifted towards a machine learning approach to better
capture the nuanced features distinguishing chirps from noise, which will be
discussed in the following section.

\subsection{Sliding convolutional classifier for chirp detection in spectrograms}
\label{sec:cnndetector}

I next employed a sliding \acrfull{CNN} to identify chirps, leveraging their
visibility in spectrograms (as shown in figure \ref{fig:multifish_chirps}).
This method involved classifying spectrogram snippets as \enquote{containing a
chirp} or \enquote{not containing a chirp} using a CNN. That is a neural
network adept at processing images through convolutional layers that extract
relevant features for classification. The process entailed moving a sliding
window across the spectrogram to generate snippets, which the CNN then
evaluated by assigning a score indicating the likelihood of a chirp's presence.
Scores were low-pass filtered and thresholded to identify chirps. The CNN was
trained on simulated chirps—simple Gaussian modulations in the frequency domain
with added noise. This model was then applied to detect chirps in real
multi fish recordings (this code is available on GitHub at:
\url{https://github.com/weygoldt/cnn-chirpdetector}).

However, this approach had limitations as well. While this method improved upon
dynamic filtering in detecting small and short chirps, it failed to reliably
identify and assign large and long chirps, which are of particular interest due
to their sparse occurrence and potential significance in specific contexts.
Additionally, the CNN struggled to differentiate between large chirps and
certain noise artifacts, and assigning chirps to the correct fish was
unreliable due to the fixed spatial context of the sliding window. In
conclusion, the sliding CNN marked progress in comparison to the dynamic
filtering approach, but it fell short in automatically detecting chirps in
complex multi-fish recordings. The following sections will detail a more
sophisticated neural network pipeline developed to address these challenges.

\vfill

\section{What is an artificial neural network?}
\label{sec:nns}

Before introducing our final approach, it's crucial to grasp the fundamentals
of a \acrfull{NN}, a mathematical model inspired by the structure of brains.
NNs are a subset of supervised learning within machine learning, designed to
learn complex non-linear relationships between input and output data, such as
classifying images into categories like \enquote{cat} or \enquote{dog}.

A NN consists of interconnected nodes or neurons arranged in layers: An input
layer receiving data, hidden layers processing the data, and an output layer
delivering the model's prediction. For instance, to classify a $(500 \times
500)$ pixel image, the input layer would have $(500 \times 500 \times 3)$
neurons for each pixel and color channel, while the output layer might have two
neurons representing the categories \enquote{cat} and \enquote{dog}. Neurons
across layers are linked by weights, the learnable parameters of the model.
Training an NN involves adjusting these weights to minimize the difference
between the model's predictions and the actual labels, a process facilitated by
an algorithm called backpropagation. This algorithm calculates the gradient of
the loss function (a measure of prediction error) with respect to the weights
and updates the weights to reduce the loss.

The simplest model, the \acrfull{MLP}, comprises multiple fully connected
layers as illustrated in figure \ref{fig:mlp}. However, MLPs alone often
struggle with complex real-world data, either failing to learn relevant
features or overfitting. This challenge is addressed by incorporating
specialized layers for different data types, such as \acrfull{LSTM} blocks for
time-dependent data or \acrshort{CNN}s for images. CNNs, for example, use
two-dimensional convolutions to extract spatial features from images, mimicking
the human visual system's ability to recognize edges, shapes, and textures.

\begin{figure}[!ht]
  \centering
  \includegraphics[width=\textwidth]{../figures/mlp.pdf}
  \caption{
    Schematic representation of a MLP with an eight-dimensional input and a
    two-dimensional output with three hidden layers. Each node is a neuron that
    is connected to all other neurons in the next layer. The thickness of the
    connecting lines corresponds to the weights (that are randomly initialized
    here). As a result, this network could be used to classify an eight
    dimensional input into two classes. Visualization by \textcite{NNSVG}.
  }
  \label{fig:mlp}
  \vspace{0.5cm}
\end{figure}

These feature extraction layers, or the \enquote{backbone} of the network, are
crucial for learning data-specific features and are followed by a
\enquote{head} responsible for making predictions. The architecture of the head
varies with the task, ranging from a single neuron for binary classification to
a softmax layer for multi-class classification or a set of neurons for object
detection, i.e., regression tasks.

In summary, NNs are powerful tools for learning from and making predictions
about complex data. Their architecture, comprising feature extraction layers
and prediction heads, can be tailored to a wide range of tasks, from text
generation to object detection. With this foundation, we can now explore how
CNNs are applied to detect chirps in recordings of multiple fish, leveraging
their ability to extract and interpret spatial features from spectrograms. The
next section will introduce the application of CNNs for chirp detection in
multi-fish recordings, showcasing the practical use of NNs in analyzing complex
biological data.

% \section{What is an Artificial Neural Network?} % (fold)
%
% Now that we already introduced the concept of a neural network, before we
% proceed to the final approach, it is important to understand the basics of
% neural networks.
%
% A neural network is a mathematical model that is inspired by the human brain.
% Neural networks fall under the category of supervised learning into the
% broader field of machine learning, meaning that they can learn a complex
% non-linear transformation between input data, for example images, and output
% data, for example the names of animals that might be in those images. It is a
% collection of connected nodes, called neurons, that are organized in layers.
% The first layer is the input layer, the last layer is the output layer and
% all layers in between are called hidden layers. The input layer is the layer
% that receives the input data, the output layer is the layer that returns the
% output of the model and the hidden layers are the layers that process the
% input data to return the output. As an example, if you would want to classify
% an image of 500 by 500 pixels with three color channels into two classes,
% let's say a \enquote{cat} or a \enquote{dog}, the input layer would consist
% of $(500 \times 500 \times 3)$ neurons and the output layer would consist of
% two, one for \enquote{cat} and one for \enquote{dog}. Classes are usually
% encoded by integer numbers, lets say 0 for \enpquote{dog} and 1 for
% \enquote{cat}.
%
% The neurons in each layer are connected to the neurons in the next layer. The
% connections between the neurons are called weights. The weights are the
% parameters of the model that are learned during training. The model is
% trained by adjusting the weights so that the output of the model is as close
% as possible to the true output. So neuron could be expressed as a function
% that takes the weighted sum of the input data and applies an activation
% function to the weighted sum. This could be just a linear function or
% logistic function. In most cases, the activation function is a non-linear
% function, such as the rectified linear unit (ReLU) \notepw{glossary} or the
% sigmoid function. This is inspired from real neurons, that (in a simplified
% way) integrate the weighted information that they receive through their
% dendrites and then, when this weighted sum is high enough to cross a
% threshold, produce a spike. The activation function is equivalent to the
% threshold crossing.
%
% The weights of each neuron are adjusted using a process called
% backpropagation. Backpropagation is the algorithm that computes the gradient
% (the derivative) of the loss function with respect to the weights of the
% model. The gradient is then used to adjust the weights so that the loss
% function is minimized. The loss function is a function that measures the
% difference between the output of the model and the true output. The loss
% function is minimized by adjusting the weights of the model. If we stick to
% our example, during each step of the training loop, the model would consume
% an image. If the image is that of a dog and the models prediction is a
% probability of 20\% for the \enquote{dog} class and 60\% for the
% \enquote{cat} class, the loss function quantifies this error from the ground
% truth (which would be 100\% and 0\% respectively). Backpropagation then
% adjusts the weights of each neuron into a direction so that the model output
% comes closer to the ground truth. And this process is repeated many times.
%
% The simplest type of neural network is the multilayer perceptron (MLP)
% \notepw{glossary}. An MLP is a type of neural network that is composed of
% multiple layers of fully connected neurons. A fully connected neuron is a
% neuron that is connected to all neurons in the previous layer. For most real
% world tasks, simply adding fully connected hidden layers is not sufficient to
% learn a good model. The model will either not be able to learn the features
% that are important to predict the output or severely overfit to the training
% data. This is were different types of feature extraction layers come into
% play. For example in the case of time-dependent data such as climate data,
% time series data or audio data, networks with a capability to remember and to
% view current events in the context of previous events are needed
% (Long-short-term memory blocks or transformer blocks such as in GPT). In case
% of images, networks with a capability to extract spatial features are needed.
% Convolutional neural networks are a type of neural network that is
% specifically designed to process images. They achieve this by including
% multiple layers of 2D convolutions between the input and output layers. The
% convolutional layers are able to extract features from the input that are
% important to classify the input into the desired classes. This is similar to
% how the human visual system works as
% \textcite{hubelReceptiveFieldsSingle1959} discovered. They found that the
% visual cortex of cats contains neurons that are specifically tuned to detect
% edges, corners and other simple features. These neurons are then connected to
% other neurons that are specifically tuned to detect more complex features
% such as shapes and textures. Convolutional neural networks are able to learn
% these features from the input data and use them to classify the input into
% the desired classes.
%
%
% These feature extraction layers are meticulously designed (by humans or by
% other neural networks that learn to optimize the architecture of neural
% networks) to be able to learn to extract features specific to the kind of
% input data. This is the so-called \enquote{backbone} of the network. The
% backbone is the part of the network that is responsible for extracting the
% features from the input data. The backbone is then connected to a head that
% is responsible for the prediction of the output. Depending on the task, the
% head can be of different architectures. For example, in the case of a binary
% classification task, the head can be a single neuron that returns the
% probability of the input data to belong to one of the two classes. In the
% case of a multi-class classification task, the head can be a
% \enquote{softmax} layer that returns the probability of the input data to
% belong to one of the classes. In the case of a detection task, the head can
% be a set of neurons that return the coordinates of the bounding boxes of the
% detected objects and the probability of the input data to belong to one of
% the classes. In reality, the model architectures currently in use are much
% more complex than the fundamental building blocks introduced here with many
% different feature extration blocks and multipurpose heads that can compute
% the loss using multiple loss functions for multiple tasks. 
%
% With this, in mind, we can now proceed to the final approach that I took to
% use machine learning to detect chirps in recordings of multiple fish. The
% next section, we will introduce how convolutional neural networks can be used
% to detect chirps in recordings of multiple fish.

% section  Convolutional Neural Networks (end) \section{The Chirp Detection
% Pipeline} % (fold) \label{sec:the_chirp_detection_pipeline}
%
% In the following sections, we will dive into the chirp detection pipeline.
% The first model, the detector, is a convolutional neural network that locates
% chirps on spectrograms. It consists of a convolutional backbone and a head
% that not only predicts the class of the object but also the coordinates of
% the bounding box for each chirp. The second model, another much simpler
% convolutional neural network, is the classifier, which assigns the detected
% chirps to the emitting fish. The two models are embedded into a pipeline
% that, when run, automatically detects and assigns chirps in recordings of
% multiple fish. Broadly, the pipeline is divided into two stages: the
% detection stage and the assignment. Both stage use different neural network
% models to process the input data and each stage has its own feature
% extraction pipeline, inference step and post-processing step. The assignment
% stage builds on the output of the detection stage and provides the final
% output of the pipeline. The pipeline takes unstructured electrode grid
% recordings as input and returns a structured output that contains the
% detected chirps and the emitting fish as well as various parameters, such as
% chirp height and duration than can be used to further analyze the detected
% chirps.
%
% Before I we introduce the full pipeline, we will first discuss the datasets
% that I used, how I generated training data for each of the models, how I
% chose the architecture of the models and how I trained the models.
%
% section The Chirp Detection Pipeline (end) chapter Background (end)
